\documentclass[11pt]{article}
\usepackage[margin=1in]{geometry}
\usepackage{amsmath,amssymb,amsthm}
\usepackage{enumitem}
\usepackage{hyperref}

\newtheorem{theorem}{Theorem}
\newtheorem{lemma}{Lemma}
\newtheorem{corollary}{Corollary}

\newcommand{\N}{\mathbb N}
\newcommand{\R}{\mathbb R}
\newcommand{\dens}{\operatorname{dens}}
\newcommand{\supp}{\operatorname{supp}}
\newcommand{\otpi}{\widehat{\otimes}_{\pi}}
\newcommand{\weakstar}{w^*}

\title{A \(c_0(\Gamma)\) Subcase of the WLD Projective Tensor Question}
\author{Run \texttt{fa\_banach\_001}, agent \texttt{agent\_lane\_10}}
\date{2026-06-18}

\begin{document}
\maketitle

\section*{Status}

This is a candidate partial result for Question 8(b) of Aviles,
Martinez-Cervantes, Rodriguez, and Rueda Zoca, \emph{Topological properties in
tensor products of Banach spaces}, arXiv:2202.00371.

\section*{Source Question}

Question 8(b) asks:

\begin{quote}
Suppose that \(X\) and \(Y\) are WLD and that either \(X\) or \(Y\) has the
Dunford-Pettis property. Is \(X\widehat{\otimes}_{\pi}Y\) WLD?
\end{quote}

The source paper proves the characterization
\[
X\otpi Y\text{ is WLD}
\]
if and only if \(X\) and \(Y\) are WLD and every operator
\[
X\to Y^*, \qquad Y\to X^*
\]
has norm-separable range.

\section*{Result}

\begin{theorem}
Let \(\Gamma\) be any index set and let \(Y\) be a WLD Banach space. Then
\[
c_0(\Gamma)\otpi Y
\]
is WLD. In particular, Question 8(b) has a positive answer whenever one factor
is \(c_0(\Gamma)\).
\end{theorem}

\section*{Proof Intuition}

The source characterization reduces the tensor question to two separable-range
operator checks. The hard-looking direction is controlled by a simple WLD
obstruction: a WLD space cannot have a nonseparable image inside any
\(\ell_1(\Gamma)\). Otherwise the image closure would be a nonseparable WLD
subspace of \(\ell_1(\Gamma)\), hence would contain \(\ell_1(\omega_1)\), but
\(\ell_1(\omega_1)\) fails Corson's property (C), while WLD spaces have
property (C) hereditarily.

For operators \(c_0(\Gamma)\to Y^*\), adjoints first give separability of the
coordinate functionals in \(\ell_1(\Gamma)\). This makes the weak-star topology
on the image ball countably determined; the source paper's WLD lemma then
upgrades weak-star separability of the image to norm separability.

\section*{Proof}

\begin{lemma}\label{lem:l1-subspace}
Every nonseparable subspace of \(\ell_1(\Gamma)\) contains an isomorphic copy
of \(\ell_1(\omega_1)\).
\end{lemma}

\begin{proof}
This is standard, but we recall the argument. Let \(Z\subseteq \ell_1(\Gamma)\)
be nonseparable. There are \(\varepsilon>0\) and an uncountable family
\((z_\alpha)\subset B_Z\) with
\(\|z_\alpha-z_\beta\|>\varepsilon\) for \(\alpha\ne\beta\). Approximate each
\(z_\alpha\) within \(\varepsilon/16\) by a finitely supported vector
\(p_\alpha\). By the finite \(\Delta\)-system lemma and a further thinning,
the supports of the \(p_\alpha\)'s have a common finite root \(R\), the root
parts \(p_\alpha|_R\) are pairwise within \(\varepsilon/16\), and the off-root
supports are pairwise disjoint.

Let \(q_\alpha=p_\alpha|_{\Gamma\setminus R}\). After discarding at most
countably many terms and thinning again, \(\|q_\alpha\|\ge \delta\) for some
\(\delta>0\). Choose the original approximation error small enough compared
with \(\delta\). For any finite scalar family \((a_\alpha)\), projection onto
\(\Gamma\setminus R\) gives
\[
\Big\|\sum_\alpha a_\alpha z_\alpha\Big\|
\ge
\Big\|\sum_\alpha a_\alpha q_\alpha\Big\|
-
\sum_\alpha |a_\alpha|\,\|z_\alpha-p_\alpha\|
\ge c\sum_\alpha |a_\alpha|
\]
for some \(c>0\), because the \(q_\alpha\)'s have disjoint supports and norms
bounded below. The reverse \(\ell_1\)-estimate is immediate from
\(\|z_\alpha\|\le 1\). Thus an uncountable subfamily of the \(z_\alpha\)'s is
equivalent to the canonical basis of \(\ell_1(\omega_1)\).
\end{proof}

\begin{lemma}\label{lem:wld-to-l1}
Let \(X\) be WLD and let \(T:X\to \ell_1(\Gamma)\) be bounded. Then \(T(X)\)
is norm separable.
\end{lemma}

\begin{proof}
Let \(Z=\overline{T(X)}\). The operator \(T:X\to Z\) has dense range, so
\(T^*:Z^*\to X^*\) is injective and weak-star to weak-star continuous. Since
\(X\) is WLD, \(X^*\) admits an injective weak-star-to-pointwise map into some
\(\ell_\infty^c(I)\). Composing with \(T^*\) shows that \(Z\) is WLD.

If \(Z\) were nonseparable, Lemma~\ref{lem:l1-subspace} would give a subspace
of \(Z\) isomorphic to \(\ell_1(\omega_1)\). But WLD spaces are hereditary and
have Corson's property (C), while \(\ell_1(\omega_1)\) fails property (C)
(as also used in Section 3 of the source paper). This contradiction proves
that \(Z\), and hence \(T(X)\), is separable.
\end{proof}

\begin{proof}[Proof of the theorem]
The space \(c_0(\Gamma)\) is WLD and has the Dunford-Pettis property. By the
source paper's Theorem 4.1, it remains to prove that every operator in both
cross directions
\[
c_0(\Gamma)\to Y^*,\qquad Y\to c_0(\Gamma)^*=\ell_1(\Gamma)
\]
has separable range.

The second direction is exactly Lemma~\ref{lem:wld-to-l1}.

For the first direction, let \(S:c_0(\Gamma)\to Y^*\) be bounded. Consider
the adjoint restricted to the canonical copy of \(Y\) in \(Y^{**}\):
\[
S^*|_Y:Y\to c_0(\Gamma)^*=\ell_1(\Gamma).
\]
By Lemma~\ref{lem:wld-to-l1}, the range \(S^*(Y)\) is separable. Choose a
countable set \((y_n)\subset Y\) such that \(S^*y_n\) is norm dense in
\(S^*(Y)\).

For \(x\in B_{c_0(\Gamma)}\) and \(y\in Y\),
\[
Sx(y)=S^*y(x).
\]
Since every \(S^*y\) is uniformly approximable on \(B_{c_0(\Gamma)}\) by some
\(S^*y_n\), the weak-star topology on \(S(B_{c_0(\Gamma)})\subset Y^*\) is
determined by the countable family \((y_n)\). Therefore the weak-star closure
of \(S(B_{c_0(\Gamma)})\) is compact metrizable, hence
\(S(B_{c_0(\Gamma)})\) is weak-star separable.

The source paper's Lemma 4.2 says that if \(E\) and \(F\) are WLD and
\(R:E\to F^*\) is bounded, then \(R(E)\) is norm separable whenever it is
weak-star separable. Applying this to \(E=c_0(\Gamma)\), \(F=Y\), and \(R=S\),
we conclude that \(S\) has norm-separable range.

Thus all cross operators have separable range, and the source characterization
implies \(c_0(\Gamma)\otpi Y\) is WLD.
\end{proof}

\section*{Verification Notes}

The proof uses only standard permanence facts for WLD spaces, the standard
subspace structure of \(\ell_1(\Gamma)\), and two explicit results from the
source paper: Theorem 4.1, which characterizes WLD projective tensor products
by separable cross-operator ranges, and Lemma 4.2, which upgrades weak-star
separability to norm separability for operators \(X\to Y^*\) when both spaces
are WLD.

\section*{Limitations and Novelty Check}

This is a partial result. It does not solve Question 8(b) for arbitrary WLD
spaces with the Dunford-Pettis property.

The run's existing attempt for arXiv:2202.00371 had identified this
\(\ell_1(\Gamma)\)-range lemma as the most promising next route, but did not
verify it. Local index searches found no existing packet for this subcase.
Web searches for the exact source question and nearby phrases found the source
arXiv page but no later direct answer.

\section*{References}

\begin{enumerate}[label={[\arabic*]}]
\item A. Aviles, G. Martinez-Cervantes, J. Rodriguez, A. Rueda Zoca,
\emph{Topological properties in tensor products of Banach spaces},
arXiv:2202.00371.
\item M. Fabian, P. Habala, P. Hajek, V. Montesinos, J. Pelant, V. Zizler,
\emph{Functional Analysis and Infinite-Dimensional Geometry}, Springer, 2001.
\end{enumerate}

\section*{Human Review Recommendation}

Review as a likely valid partial result. The proof point most worth checking is
Lemma~\ref{lem:l1-subspace}, the standard extraction of an
\(\ell_1(\omega_1)\)-sequence from a nonseparable subspace of
\(\ell_1(\Gamma)\).

\end{document}
